\chapter{Rotating the net}\label{ch:rotation}
\begin{goals}
\textbf{In this chapter} you turn a net by dragging it or by naming a direction for its center, and learn
what a rotation changes (the picture) and what it doesn't (the measurements, the selections, the
statistics).
\end{goals}

A net normally looks straight down: the center is vertical, north at the top. Turning it brings another
direction to the center, which is how you look along a fold axis, straighten a girdle into a great circle
through the middle, or check that a cluster near the primitive is one cluster and not two halves.

\section{Dragging}\label{sec:rotate-drag}

With the \ui{pick} tool chosen (\cref{sec:select-tools}), press on the net and drag. The net turns as if
you were rolling a ball under your hand: the point you pressed follows the pointer, and the cardinal
letters, the graticule, the contours and every layer turn with it. A press that moves less than a few
pixels is still a click, so picking works as before; the click that ends a drag doesn't pick.

The other tools use a drag for their shapes, so switch back to \ui{pick} to turn the net.

\begin{why}
wuffle turns the net with an \term{arcball}: each small movement of the pointer adds a small rotation to
the one the net already has. Unlike turning by two angles (a trend and a plunge for the center), it has no
direction where the drag suddenly flips or stops responding, so you can roll the net as far as you like in
any direction.
\end{why}

The rotation belongs to the diagram. It is kept with the project, every view of every layer on that net
shares it, and a figure made from the net (\cref{ch:figures}) is drawn as rotated.

\section{Naming the center}\label{sec:rotate-centre}

The console's \code{rotate} brings a direction to the center of the active net:\index{rotate (command)}

\begin{codeblock}
rotate 120→40       # Net 1: centered on 120→40
rotate '120>40'     # the same, typed without the arrow (quoted: > is the shell's)
rotate off          # Net 1: not rotated
\end{codeblock}

The direction is a line. \code{rotate 045/30} refuses with \ui{rotate 045→30 (the direction to put at the
center)}; to look down a plane's pole, give the pole as a line. Only a net rotates: with a rose or a
fabric plot active, \code{rotate} says \ui{only a net rotates}.

\shot{rotated-net}{The Tocher net centered on \code{120→40} (the cardinal letters have moved with it) and
the same cone selection as \cref{fig:select-cone}, the same 18 poles, drawn where they now fall.}

\section{Undoing and resetting}\label{sec:rotate-reset}

A drag is one undo step, \ui{Undo diagram settings}, however long you roll the net, as long as you keep
moving. A pause of most of a second starts a new step, so a drag in two movements can be undone one movement
at a time. \code{rotate} is one step too; and as with a drag, \code{rotate} lines run within a moment
of each other on one net (as in a routine) merge into one step.

To go back to the plain net, right-click the net in the tree and choose \ui{Reset rotation}\index{Reset
rotation} (grayed out when the net isn't rotated), or type \code{rotate off}. Either clears a dragged
rotation or a named center.

\section{What a rotation doesn't change}\label{sec:rotation-true}

A rotation changes the picture only. Everything wuffle computes is in geographic directions:

\begin{itemize}
\item \textbf{Selections.} A shape is on the sphere (\cref{sec:select-sphere}), so a cone drawn on a
rotated net is a cone about the same direction as on the plain one. \code{select cone 040/50 20} selects
the same 18 Tocher poles whether the net is plain, centered on \code{120→40} or rolled by hand.
\item \textbf{Picking and hovering.} The status line reads the direction under the pointer in the real
world, not in the net's turned frame. Hover a pole and it reads the attitude in the table.
\item \textbf{Statistics and derived data.} The card, the mean, the eigenvectors and the $\beta$ diagram
describe the measurements, which don't move.
\item \textbf{The data.} Nothing in the tables changes.
\end{itemize}

\section{In the files}

A rotation is one of the diagram's settings, under \code{"props"} in \file{diagrams/<name>/diagram.json}
(\cref{ch:project-fs}). A named center is stored as \code{"rotation": [120, 40]}; a dragged rotation as
\code{"matrix"}, the nine numbers of the rotation matrix. Writing \code{"rotation": null} and
\code{"matrix": null} there resets the net, as \ui{Reset rotation} does. The net's \file{net.svg} is drawn
as rotated.

\begin{tryit}
\begin{enumerate}
\item Open the Tocher sample, type \code{contour kamb 3}, then \code{select cone 040/50 20}. Note where the
18 violet poles are.
\item Make sure \ui{pick} is chosen, and drag the net slowly from its center toward the east. The net
rolls; the violet poles move with everything else and stay violet.
\item Type \code{select cone 040/50 20} again. It still selects 18.
\item Type \code{rotate 113→48}, Tocher's V3 (\cref{sec:stats-eigen}). You are looking down the axis of
the girdle, which now runs around the primitive.
\item Type \code{rotate off}.
\end{enumerate}
\end{tryit}
